Un stock de \np{1000} graines comporte des graines de basilic et des graines d'estragon. Une fois semées, certaines graines germent, d'autres ne germent pas. On donne ci-dessous la répartition des \np{1000} graines.

\begin{center}
\begin{tabular}{|l|c|c|c|}
	\cline
	{ 2 - 4 }
\multicolumn{1}{c|}{} & Graines de basilic & Graines d'estragon &   Total   \\
	\hline
	Graines qui germent             & $x$ & 100 & 250 \\
	\hline
	Graines qui ne germent pas      & 450 & $y$ & 750  \\
	\hline
	Total                           & 600 & 400 & \np{1000} \\
	\hline
\end{tabular}
\end{center}

\begin{enumerate}
\item Déterminer les valeurs de $x$ et $y$.

Interpréter ces résultats dans le contexte de l'exercice.
\item On sème au hasard une graine de ce stock.

On note les évènements :

\begin{itemize}[label={\textbullet~}]
	\item $B$	: \og la graine est une graine de basilic \fg{};
	\item $E$ : \og la graine est une graine d'estragon \fg{};
	\item $G$ : \og la graine germe \fg{};
	\item $\overline{G}$ : l'évènement contraire de $G$.
\end{itemize}

\begin{enumerate}
	\item Déterminer les probabilités $P(B)$ et $P(E)$.

	\item Définir à l'aide d'une phrase l'évènement $E \cap G$, puis calculer sa probabilité.

	\item Sachant que la graine semée a germé, quelle est la probabilité que ce soit une graine d'estragon?

	\item Sachant que la graine semée est une graine d'estragon, quelle est la probabilité qu'elle germe ?

	\item Un jardinier affirme : \og Le fait qu'une graine germe ne dépend pas du type de graine semée (basilic ou estragon). \fg

	Cette affirmation est-elle vraie?

	Justifier la réponse.

	\end{enumerate}
\end{enumerate}
